% -----------------------------------------------------------
\section{Gospel}
% -----------------------------------------------------------
	\label{GOSPEL}
	
% % ----------------------
% \subsection{See also}
% % ----------------------

% % ----------------------
% \subsection{1828 American Dictionary of the English Language} % *1828
% % ----------------------
	% \label{GOSPEL-1828}

	% \begin{compactitem}
		% \item
	% \end{compactitem}

% % ----------------------
% \subsection{Oxford English Dictionary} % *OED
% % ----------------------
	% \label{GOSPEL-OED}

	% \begin{compactitem}
		% \item[]
	% \end{compactitem}

% ----------------------
\subsection{Scriptural Usage} % *SCRIPTURES
% ----------------------
	\label{GOSPEL-SCRIPTURES}

	% % -----
	% \subsubsection{Old Testament} % *OT
	% % -----
		% \label{GOSPEL-OT}
	
		% % --
		% \paragraph{KJV}
		% % --
			% \label{GOSPEL-OT-KJV}
	
			% \begin{tabular}{ll}
				% Total occurrences in the OT & 
			% \end{tabular}
			
			% \begin{compactdesc}
				% \item[]
			% \end{compactdesc}
			
		% % --
		% \paragraph{Hebrew Bible}
		% % --
			% \label{GOSPEL-HEBREW}
		
			% %
			% \subparagraph{\texorpdfstring{\he{sample word}}{}} % paste actual hebrew text here
			% %
				% \label{PAGA} % whatever is in step bible without periods
			
				% \begin{compactdesc}
					% \item[HALOT , PDF ] \textbf{} % Use the 1994 PDF
						% \begin{compactitem}
							% \item[1]
						% \end{compactitem}
					% \item[BDB , PDF ] \textbf{}
						% \begin{compactitem}
							% \item[1]
						% \end{compactitem}
					% \item[DCH PDF ] \textbf{}
						% \begin{compactitem}
							% \item[1]
						% \end{compactitem}
				% \end{compactdesc}
				
				% \begin{compactdesc}
					% \item[Some OT uses] \textbf{}
						% \begin{compactdesc}
							% \item[]
						% \end{compactdesc}
				% \end{compactdesc}
			
		% % --
		% \paragraph{Septuagint}
		% % --
			% \label{GOSPEL-LXX}
		
			% %
			% \subparagraph{\texorpdfstring{\gr{sample word}}{}}
			% %
		
	% -----
	\subsubsection{New Testament} % *NT
	% -----
		\label{GOSPEL-NT}
	
		% --
		\paragraph{KJV}
		% --
			\label{GOSPEL-NT-KJV}
			
	
			\begin{tabular}{ll}
				Total occurrences in the NT & 72
					% From the step bible, a search for εὐαγγέλιον, it appears in 72 verses and those verses take 72 lines when you copy and paste them, so no verse has the word twice I think
			\end{tabular}
			
			\begin{compactdesc}
				\item[{\hyperref[ROMANS-1-16]{Romans 1:16}}]
			\end{compactdesc}
			
		% --
		\paragraph{Greek New Testament} % *GREEK
		% --
			\label{GOSPEL-GREEK}
		
			%
			\subparagraph{\texorpdfstring{\gr{εὐαγγέλιον}}{}}
			%
				\label{EUAGGELION}
			
				\begin{compactdesc}
					\item[BDAG 355a] originally `a reward for good news', then simply `good news'
						\begin{compactitem}
							\item[1] \textbf{God's good news to humans, \textit{good news}} as proclamation
							\item[1.a] absolute
							\item[1.a.$\alpha$] \gr{τὸ εὐαγγέλιον}
							\item[1.a.$\beta$] in gen., dependent on another noun \gr{ὁ λόγος τοῦ εὐ.}...\gr{τὸ μυστήριον τ. εὐ.}...\gr{ἡ ἀλήθεια τοῦ εὐ.}...etc.
							\item[1.a.$\gamma$] in certain combinations with verbs \gr{τὸ εὐ. κηρύσσειν}...\gr{καταγγέλλειν}...etc.
							\item[1.b] in combination
							\item[1.b.$\alpha$] with adjective
							\item[1.b.$\beta$] with genitive
							\item[1.b.$\beta$.\heb{'}] objective genitive
							\item[1.b.$\beta$.\heb{b}] subjective genitive
							\item[2] \textbf{details relating to the life and ministry of Jesus, \textit{good news of Jesus}}
							\item[3] \textbf{a book dealing with the life and teaching of Jesus, \textit{a gospel account}} that deals with the life and teaching of Jesus
						\end{compactitem}
					% \item[CGL , PDF ] \textbf{}
						% \begin{compactitem}
							% \item 
						% \end{compactitem}
					% \item[LSJ , PDF ] \textbf{}
						% \begin{compactitem}
							% \item 
						% \end{compactitem}
				\end{compactdesc}
				
				% \begin{compactdesc}
					% \item[Some NT uses] \textbf{}
						% \begin{compactdesc}
							% \item[]
						% \end{compactdesc}
				% \end{compactdesc}
		
	% % -----
	% \subsubsection{Book of Mormon} % *BOM
	% % -----
		% \label{GOSPEL-BOM}
	
		% \begin{tabular}{ll}
			% Total occurrences in the BOM & 
		% \end{tabular}
		
		% \begin{compactdesc}
			% \item[]
		% \end{compactdesc}
		
	% % -----
	% \subsubsection{Doctrine and Covenants} % *DC
	% % -----
		% \label{GOSPEL-DC}
	
		% \begin{tabular}{ll}
			% Total occurrences in the DC & 
		% \end{tabular}
		
		% \begin{compactdesc}
			% \item[]
		% \end{compactdesc}
		
	% % -----
	% \subsubsection{Pearl of Great Price} % *PGP
	% % -----
		% \label{GOSPEL-PGP}
	
		% \begin{tabular}{ll}
			% Total occurrences in the PGP & 
		% \end{tabular}
		
		% \begin{compactdesc}
			% \item[]
		% \end{compactdesc}
		
% ----------------------
\subsection{Key Phrases} % *KEYPHRASES *PHRASES
% ----------------------
	\label{GOSPEL-PHRASES}

	\begin{compactdesc}
		\item[fulness of * gospel] 
		\item[fulness of my gospel] 
		\item[fulness of the gospel] 
		\item[law of the gospel] See especially the \hyperref[TEMPLE-ENDOWMENT-ENDOW-PRE-1990-GOSPEL]{endowment}. \textbf{\#times} | list
			% \begin{compactdesc}
				% \item[Bible anchor verses] \phantom{.}
					% \begin{compactdesc}
						% \item[Romans 5:5] \gr{}
					% \end{compactdesc}
				% \item[Commentary] ---
			% \end{compactdesc}
	\end{compactdesc}
	
% % ----------------------
% \subsection{Bible Dictionaries} % *DICTIONARIES
% % ----------------------
	% \label{GOSPEL-BD}

% % ----------------------
% \subsection{Restoration Usage} % *RESTORATION
% % ----------------------
	% \label{GOSPEL-RESTORATION}

	% % -----
	% \subsubsection{Joseph Smith} % *JS
	% % -----
		% \label{GOSPEL-JS}
	
		% \begin{compactdesc}
			% \item[{\hyperref[KEY]{Date/Source}}]
		% \end{compactdesc}
		
	% % -----
	% \subsubsection{Temple Ordinances}
	% % -----
		% \label{GOSPEL-TEMPLE}
	
		% \begin{compactdesc}
			% \item[{\hyperref[KEY]{Date/Source}}]
		% \end{compactdesc}
	
	% % -----
	% \subsubsection{Brigham Young} % *BY
	% % -----
		% \label{GOSPEL-BY}
	
		% \begin{compactdesc}
			% \item[{\hyperref[KEY]{Date/Source}}]
		% \end{compactdesc}
	
% ----------------------
\subsection{Commentary and Studies} % *COMMENTARY *STUDIES
% ----------------------
	\label{GOSPEL-COMMENTARY}

	% % -----
	% \subsubsection{Key Points}
	% % -----
		% \label{GOSPEL-KEYS}
	
		% \begin{compactitem}
			% \item
		% \end{compactitem}
		
	% -----
	\subsubsection{Scriptural definitions of the gospel}
	% -----
		\label{GOSPEL-definition}
		
		\begin{compactdesc}
			\item[{\hyperref[ROMANS-1-16]{Romans 1:16}}]
			\item[{\hyperref[3NEPHI-27-13]{3 Nephi 27:13--15}}]
			\item[{\hyperref[3NEPHI-27-16]{3 Nephi 27:16--21}}] \phantom{.}
				\begin{compactitem}
					\item This really goes with the one before, but 16--21 tell us more about the gsopel and then in 21 he says, ``This is my gospel,'' and it seems like he's talking about the things he just said.
				\end{compactitem}
			\item[{\hyperref[DC6-29]{DC 6:29}}] \phantom{.}
				\begin{compactitem}
					\item Notice how this verse refers to ``this part of my gospel.'' So it's possible to have one part and not have another, and it's possible for some parts to be held back because of wickedness, which is what verse 26 says. Note that that can be personal wickedness or group wickedness.
				\end{compactitem}
			\item[{\hyperref[DC10-67]{DC 10:67}}] 
			\item[{\hyperref[DC11-24]{DC 11:24}}]
			\item[{\hyperref[DC33-12]{DC 33:12}}] \phantom{.}
				\begin{compactitem}
					\item See verses 10--13 here.
				\end{compactitem}
			\item[{\hyperref[DC39-6]{DC 39:6}}] 
			\item[{\hyperref[DC42-12]{DC 42:12}}] 
			\item[{\hyperref[DC66-2]{DC 66:2}}]
			\item[{\hyperref[DC76-13]{DC 76:13--14}}]
			\item[{\hyperref[DC76-40]{DC 76:40--42}}]
			\item[{\hyperref[DC79-1]{DC 79:1}}]
		\end{compactdesc} 